\section{Engram Protocol}
\label{sec:protocol}

The mechanisms follow the order of the data: the architecture, the sensors that
produce observations, the certificate that aggregates them, the state machine
that consumes the aggregate, the validity predicate that enforces it, and the
recovery path out of \textsc{Sovereign}.

\subsection{Architecture Overview}
\label{subsec:architecture}

Engram is a Cosmos SDK application~\cite{kwon2016cosmos} running on a
CometBFT core. Three peripheral layers are monitored: (i) Bitcoin, used as a
settlement and finality anchor via periodic checkpoints; (ii)
Celestia~\cite{albassam2019lazyledger}, used as the DA layer; and (iii) the
peer-to-peer interfaces through which validators reach the Engram, Bitcoin,
and Celestia networks. Engram follows one rule throughout:
\emph{sensors observe, the quorum certifies, the state machine decides}.
Each validator's sensors read the peripheral layers
locally and asynchronously, so two honest validators may observe different
values at the same instant. Those readings never enter replicated state
directly. Instead, each validator signs its reading as a CometBFT vote
extension; the next proposer aggregates a quorum of them into a
\emph{health certificate} (Section~\ref{subsec:health-certificate}); and the
FSM transition is a deterministic function of that certificate and the
previously committed FSM state (Fig.~\ref{fig:architecture}).

\begin{figure}[h]
\centering
\definecolor{zoneLocalFill}{RGB}{243,244,246}
\definecolor{zoneLocalLine}{RGB}{107,114,128}
\definecolor{zoneRepFill}{RGB}{238,240,252}
\definecolor{zoneRepLine}{RGB}{67,56,202}
\resizebox{\linewidth}{!}{%
\begin{tikzpicture}[
    >=Stealth, font=\footnotesize,
    box/.style={draw, rounded corners=3pt, align=center, minimum height=0.9cm, inner sep=4pt, thick},
    ext/.style={box, fill=white, draw=zoneLocalLine, minimum width=2.8cm},
    loc/.style={box, fill=white, draw=zoneLocalLine, text width=11.4cm},
    det/.style={box, fill=white, draw=zoneRepLine, text width=11.4cm},
    flow/.style={->, thick, black!75},
    lbl/.style={font=\scriptsize, align=center}
]
% Peripheral layers (external, not part of consensus)
\node[ext] (cel) {Celestia\\(DA)};
\node[ext, left=0.6cm of cel] (btc) {Bitcoin\\(settlement)};
\node[ext, right=0.6cm of cel] (p2p) {Peripheral P2P\\(peer quality)};
% Validator-local zone
\node[loc, below=1.2cm of cel] (sens) {Local sensors: BTC gap, anchor re-verification, DA gap, 6-bound P2P profile};
\node[loc, below=0.6cm of sens] (ve) {\texttt{ExtendVote}: signed, height- and round-bound observation $o_i$};
% Replicated, deterministic zone
\node[det, below=2.0cm of ve] (cert) {\texttt{PrepareProposal}: quorum of observations $\Rightarrow$ certificate $C^{h}$ as $\mathit{Txs}[0]$};
\node[det, below=0.6cm of cert] (pp) {\texttt{ProcessProposal}: re-verify signatures, re-run $\mathsf{Agg}$, check $\delta$, receipts, proof flag};
\node[det, below=0.6cm of pp] (fsm) {CometBFT prevote/precommit $\to$ \texttt{PreBlocker}: commit $s^{h}$, circuit breaker};
% Zones, drawn behind the boxes
\begin{scope}[on background layer]
  \draw[rounded corners=6pt, fill=zoneLocalFill, draw=zoneLocalLine, dashed, thick]
    ([xshift=-0.4cm, yshift=0.3cm]sens.north west) rectangle ([xshift=0.4cm, yshift=-0.75cm]ve.south east);
  \draw[rounded corners=6pt, fill=zoneRepFill, draw=zoneRepLine, thick]
    ([xshift=-0.4cm, yshift=0.3cm]cert.north west) rectangle ([xshift=0.4cm, yshift=-0.75cm]fsm.south east);
\end{scope}
\node[font=\scriptsize\bfseries, text=zoneLocalLine!80!black, anchor=south west]
  at ([xshift=-0.3cm, yshift=-0.7cm]ve.south west) {VALIDATOR-LOCAL \textnormal{\textperiodcentered{} live}};
\node[font=\scriptsize\bfseries, text=zoneRepLine, anchor=south west]
  at ([xshift=-0.3cm, yshift=-0.7cm]fsm.south west) {REPLICATED \textnormal{\textperiodcentered{} deterministic}};
% Data flow
\draw[flow] (btc) -- (sens.north -| btc);
\draw[flow] (cel) -- (cel |- sens.north);
\draw[flow] (p2p) -- (sens.north -| p2p);
\draw[flow] (sens) -- (ve);
\draw[flow, zoneRepLine, very thick] (ve) -- node[lbl, right, pos=0.5, fill=white, inner sep=1.5pt, text=black]{vote extensions of height $h{-}1$} (cert);
\draw[flow] (cert) -- (pp);
\draw[flow] (pp) -- (fsm);
% Recovery path (off-chain prover, outside the zones)
\node[box, dashed, fill=white, draw=black!60, minimum width=2.4cm, right=1.1cm of $(pp.east)!0.5!(fsm.east)$, anchor=west] (prover) {ZK prover\\(off-chain)};
\draw[->, dashed, thick, black!60] (fsm.east) -| node[lbl, pos=0.25, below]{headers} (prover.south);
\draw[->, dashed, thick, black!60] (prover.north) |- node[lbl, pos=0.7, above]{proof tx} (pp.east);
\end{tikzpicture}}
\caption{Engram architecture. Live sensor reads stay in the validator-local
zone and enter consensus only as signed vote extensions (bold arrow, the single
crossing). The replicated zone, from certificate aggregation downward, is a
pure function of data every validator holds, so the committed FSM state
$s^{h}$ is deterministic. The dashed path is the recovery proof
(Section~\ref{subsec:reanchoring}).}
\label{fig:architecture}
\end{figure}

\subsection{Network Sensors}
\label{subsec:sensors}

Sensors are the only non-deterministic component; their outputs enter
consensus only through the certificate (Section~\ref{subsec:health-certificate}).

\textbf{Bitcoin finality gap.} $\Delta H_{\text{BTC}} = H_{\text{current}} -
\min(H_{\text{submitted}}, H_{\text{anchored}})$ measures how far the latest
checkpoint lags Bitcoin's tip; it grows under congestion or checkpoint
censorship. Because a delay metric cannot detect a checkpoint that was
reorganized away, each validator also re-verifies, against its own Bitcoin
full node, that the checkpoint transaction is still included at depth
$\geq K_{\text{DEEP\_FINALITY}}$. A failed re-verification is a critical
condition.

\textbf{DA gap.} $\Delta H_{\text{DA}} = H_{\text{local}} -
H_{\text{verified}}$ measures the lag between the chain head and the highest
height whose DA submission the validator has confirmed as retrievable from
its own Celestia bridge node. This is bridge-level retrievability, not a
data-availability-sampling proof; Section~\ref{subsec:sim-to-real} describes
what the prototype publishes.

\textbf{P2P health.} Peer count alone cannot distinguish a long-lived honest
peer from a fresh Sybil, and eclipse attacks combine vectors no single
metric captures~\cite{rehman2025eclipse,heilman2015eclipse}. The profiler
combines six bounds -- subnet diversity, active anchor peers, clean peers,
churn rate, mean peer tenure, and latency -- applied to the Engram P2P
interface, the Bitcoin client, and the Celestia client. These bounds
characterize peripheral and partial P2P degradation; a partition of the
core consensus network below a quorum is outside what any FSM can handle
(Section~\ref{sec:intro}).

\textbf{Agreement versus correctness of health.} Lemma~\ref{lem:honest-range}
(Section~\ref{subsec:health-certificate}) bounds what
Byzantine validators can do to the aggregate; it does not make an honest
sensor accurate. The six P2P bounds are a heuristic detector: an eclipse that
keeps peer count, subnet diversity, tenure, and latency within bounds while
controlling a validator's view is not detected by them. The TLA+ model assumes
that an eclipsed honest validator reports a degraded reading and never a
healthier one
(Section~\ref{sec:formal}); that is a sensor-soundness assumption of the
model, not a property derived from the six bounds.

\subsection{Health Certificates}
\label{subsec:health-certificate}

The health certificate turns asynchronous, possibly inconsistent local
observations into a single value that every validator can recompute.

\textbf{Observation.} At height $h{-}1$ and round $r$, validator $i$ produces
an observation $o_i = (i, m_i, \sigma_i)$. Here $m_i$ is its local metric
vector (Section~\ref{subsec:sensors}) -- Bitcoin gap, anchored and current Bitcoin heights, anchor
re-verification result, DA gap and verified height, DA sampling and
attestation flags, and the six P2P-quality bounds -- and $\sigma_i$ is $i$'s
signature over CometBFT's vote-extension sign bytes
$(\mathit{chainID}, h{-}1, r, m_i)$. The extension is produced in
\texttt{ExtendVote}, the only ABCI++ hook permitted to read live sensors.

\textbf{Certificate.} Let $r^{\star}$ be the round in which height $h{-}1$
committed. The proposer of height $h$ receives the extended commit of height
$h{-}1$ and places
\[
C^{h} = \bigl(h{-}1,\; r^{\star},\; \{o_i\}_{i \in S}\bigr)
\]
as the first transaction of its proposal. $C^{h}$ is \emph{well-formed} iff
(a) every $\sigma_i$ verifies against the genesis key of validator $i$ for
$(h{-}1, r^{\star})$; (b) each validator contributes at most one observation;
(c) every $i \in S$ has a commit precommit in round $r^{\star}$ of height
$h{-}1$; and (d) $\operatorname{wt}(S) > 2W/3$, so $S$ is a quorum.

\textbf{Aggregation.} For a numeric metric $j$, the aggregate is the smallest
value whose cumulative weight exceeds half of the certificate weight; a
Boolean metric is true iff the weight backing true exceeds half of it:
\begin{align}
\widehat m_j &= \min\bigl\{x : \operatorname{wt}(\{i \in S : m_{i,j} \le x\}) > \tfrac12 \operatorname{wt}(S)\bigr\}, \label{eq:agg}\\
\widehat b_j &= \mathbf{1}\bigl[\operatorname{wt}(\{i \in S : b_{i,j} = 1\}) > \tfrac12 \operatorname{wt}(S)\bigr]. \nonumber
\end{align}
$\mathsf{Agg}(C^{h})$ collects these values. The Boolean metrics are failure
flags (failed anchor re-verification, failed DA sampling, failed DA
attestation), so a tie leaves the flag unset. $\mathsf{Agg}$ is a pure
function: every validator that receives $C^{h}$ computes the same value
without contacting any peripheral network. The TLA+ model gives every
validator the same voting power (Section~\ref{sec:formal}).

\refstepcounter{lemma}\label{lem:honest-range}\textbf{Lemma~\thelemma{} (honest-bounded aggregate).} \emph{Let $C^{h}$ be well-formed
and $\operatorname{wt}(B) < W/3$. For every numeric metric $j$,}
\[
\min_{i \in S \cap H} m_{i,j} \;\le\; \widehat m_j \;\le\; \max_{i \in S \cap H} m_{i,j},
\]
\emph{and every Boolean metric on which all members of $S \cap H$ agree takes
that value.} \emph{Proof sketch.} Since $\operatorname{wt}(S) > 2W/3$ and
$\operatorname{wt}(S \cap B) < W/3$, the Byzantine weight in $S$ is below
$\operatorname{wt}(S)/2$, so the honest weight in $S$ exceeds
$\operatorname{wt}(S)/2$. Let $a$ and $b$ be the smallest and largest honest
values. The weight of values at most $b$ is at least the honest weight, so
$\widehat m_j \le b$. The weight of values below $a$ is at most the
non-honest weight, which is below $\operatorname{wt}(S)/2$, so no value below
$a$ satisfies Eq.~\eqref{eq:agg} and $\widehat m_j \ge a$. For a Boolean
metric on which all honest members agree on $v$, the weight backing $v$
exceeds half, and the weight backing the other value is below half.
$\square$

Lemma~\ref{lem:honest-range} does not require honest observations to agree; it bounds what the
adversary can do when they disagree. Each adversarial case reduces to it:
\begin{itemize}
  \item \emph{Stale or replayed observation.} A signature binds the height
  and round; an observation signed at another height or round fails (a) and
  is discarded.
  \item \emph{Conflicting observations.} Two different signed observations
  by one validator count once by (b); Byzantine weight is below $W/3$ either
  way.
  \item \emph{False observations.} Byzantine validators below $W/3$
  reporting extreme values cannot move the aggregate outside the honest range
  (Lemma~\ref{lem:honest-range}).
  \item \emph{Proposer omission.} The proposer chooses $S$, and may omit
  honest observations. If $S$ still meets (d), Lemma~\ref{lem:honest-range} still holds. The
  proposer's residual influence is choosing \emph{which} honest value in the
  honest range is selected when honest validators genuinely disagree. If
  omission drops $S$ to weight at most $2W/3$, the certificate is not
  well-formed and honest validators reject the proposal. The only exception
  is the first height after vote extensions are enabled, whose predecessor
  carries no extensions; there the certificate is replaced by a maximally
  degraded metric vector, the conservative choice. \pending{re-run E8 live
  with the omission attack (B3).}
  \item \emph{Missing or malformed observations.} CometBFT accepts a
  precommit's extension only if its signature verifies under the validator's
  consensus key, and the application's \texttt{VerifyVoteExtension} accepts
  only extensions that decode; signatures are checked again against the
  genesis keys when the certificate is aggregated. The extended commit for
  $h{-}1$ carries weight above $2W/3$, so an honest proposer can always build
  a well-formed certificate, and a Byzantine validator cannot reduce the
  usable weight by sending a decodable extension with a bad signature,
  because CometBFT would not have accepted that vote. This relies on CometBFT
  and the aggregation using the same key and sign bytes.
  \item \emph{Divergent honest views.} When honest views straddle a
  threshold, different proposers may select certificates on different sides
  of it. Each resulting transition is legal and deterministic given its
  certificate; the FSM's hysteresis absorbs short alternations
  (Section~\ref{subsec:transitions}).
\end{itemize}

\textbf{Height and round binding.} The certificate committed at height $h$
describes observations signed at height $h{-}1$, one block of lag. A
signature binds $(\mathit{chainID}, h{-}1, r, m_i)$, and a well-formed
certificate takes observations only from the commit round $r^{\star}$ and
only from validators whose precommit is in that commit.
This prevents replay of an observation from another height or round into
$C^{h}$. The binding is a property of the consensus instance, not of the
measurement: a validator can still sign a cached Bitcoin tip or stale peer
metrics (Section~\ref{subsec:sensors}).

\subsection{Finite State Machine}
\label{subsec:fsm}

The FSM has four states (Table~\ref{tab:fsm-states}). \textsc{Anchored} is
normal operation: blocks are valid only with a Bitcoin receipt and a DA
receipt within tolerance, and withdrawals are permitted.
\textsc{Suspicious} is an early-warning state reached after a sustained
warning; withdrawals remain permitted, but admission of forced transactions
is rate-limited by a bounded queue. \textsc{Sovereign} is the fallback: the
Bitcoin and DA receipt checks are relaxed for each layer that the certificate
does not report healthy, blocks continue under the validator quorum, and the
circuit breaker locks all cross-chain withdrawals. \textsc{Recovering} is
entered once the certificate reports healthy peripherals; withdrawals stay
locked until the exit gate of Section~\ref{subsec:reanchoring} is satisfied.

\begin{table}[h]
\caption{FSM states and their validity requirements. The exit gate is defined
in Section~\ref{subsec:reanchoring}.}
\label{tab:fsm-states}
\centering
\footnotesize
\newcommand{\rcc}[1]{\raggedright\arraybackslash #1}
\begin{tabular}{@{}>{\rcc}p{0.12\linewidth}>{\rcc}p{0.31\linewidth}>{\rcc}p{0.13\linewidth}>{\rcc}p{0.15\linewidth}>{\rcc}p{0.18\linewidth}@{}}
\toprule
\textbf{State} & \textbf{Validity requires} & \textbf{Withdrawals} &
\textbf{Forced-tx admission} & \textbf{BTC settlement} \\
\midrule
\textsc{Anchored} & Quorum; BTC and DA receipts & Permitted & Full & Yes ($K_{\text{DEEP}}$ conf.) \\
\textsc{Suspicious} & Quorum; receipt of each layer reported healthy & Permitted & Rate-limited & Yes, delayed \\
\textsc{Sovereign} & Quorum; receipt of each layer reported healthy & Locked & Full & No \\
\textsc{Recovering} & Quorum; BTC and DA receipts; exit gate & Locked & Full & Pending proof \\
\bottomrule
\end{tabular}
\\[2pt]
\footnotesize CometBFT's deterministic finality is unaffected by the FSM state;
the measured commit interval on our testbed is about 1--2\,s in every state.
Bitcoin settlement is probabilistic and follows Bitcoin's own confirmation
time ($K_{\text{DEEP}}$ blocks, $\sim$10\,min per block on mainnet); it is
not measured in this paper.
\end{table}

Fig.~\ref{fig:fsm} shows the transition function $\delta$. A proposal
carries a target state $s^{h}$; honest validators prevote for it only if
\begin{equation}
s^{h} = \delta\bigl(s^{h-1}, \mathsf{Agg}(C^{h})\bigr),
\label{eq:delta}
\end{equation}
where $\delta$ also reads the hysteresis counters committed with $s^{h-1}$
(Section~\ref{subsec:transitions}). Since $\mathsf{Agg}$ and $\delta$ are
deterministic and the state and counters are committed, every honest
validator reaches the same verdict on a given proposal, and the committed
state inherits CometBFT's agreement guarantee. $\delta$ evaluates the
conditions in a fixed priority: a critical condition first, then a warning
(in \textsc{Recovering}, any non-healthy reading), then a healthy reading.
The first condition that applies decides, so a critical condition overrides
every hysteresis counter.

\begin{figure}[h]
\centering
\resizebox{\linewidth}{!}{%
\begin{tikzpicture}[
    >=Stealth, auto,
    state/.style={rectangle, rounded corners, draw, very thick, minimum width=1.9cm,
                  minimum height=1.1cm, align=center, font=\footnotesize\bfseries},
    lab/.style={font=\scriptsize, align=center, inner sep=1.5pt, text width=2.3cm}
]
\node[state, fill=yellow!35, draw=yellow!50!black] (anchored) {ANCHORED};
\node[state, fill=orange!45, draw=orange!60!black, right=2.9cm of anchored] (suspicious) {SUSPICIOUS};
\node[state, fill=red!40, draw=red!60!black, right=2.9cm of suspicious] (sovereign) {SOVEREIGN};
\node[state, fill=green!35, draw=green!50!black, right=2.9cm of sovereign] (recovering) {RECOVERING};

% Degradation above the boxes, recovery below.
\draw[->] ([yshift=0.2cm]anchored.east) -- node[lab, above]{Sustained\\warning} ([yshift=0.2cm]suspicious.west);
\draw[->] ([yshift=-0.2cm]suspicious.west) -- node[lab, below]{Sustained\\healthy} ([yshift=-0.2cm]anchored.east);
\draw[->] (suspicious.east) -- node[lab, above]{Critical\\or timeout} (sovereign.west);
\draw[->] ([yshift=0.2cm]recovering.west) -- node[lab, above, text width=2.8cm]{Critical or\\sustained unhealthy} ([yshift=0.2cm]sovereign.east);
\draw[->] ([yshift=-0.2cm]sovereign.east) -- node[lab, below]{Healthy} ([yshift=-0.2cm]recovering.west);
\draw[->] (anchored.north) to[bend left=32] node[lab, above, text width=4cm]{Critical (immediate)} (sovereign.north);
\draw[->] (recovering.south) to[bend left=32] node[lab, below, yshift=-2pt, text width=4cm]{Stable healthy + proof verified} (anchored.south);
\end{tikzpicture}}
\caption{FSM transition function $\delta$. \emph{Sustained} means the condition
accumulates over several blocks (hysteresis; up-counters tolerate isolated
non-healthy blocks); critical conditions skip the wait and act immediately.
The waits are given in Table~\ref{tab:fsm-waits}.}
\label{fig:fsm}
\end{figure}

\begin{table}[h]
\caption{Hysteresis waits of the FSM (defaults of the prototype;
Section~\ref{subsec:transitions}).}
\label{tab:fsm-waits}
\centering
\footnotesize
\newcommand{\rcc}[1]{\raggedright\arraybackslash #1}
\begin{tabular}{@{}>{\rcc}p{0.2\linewidth}>{\rcc}p{0.34\linewidth}>{\rcc}p{0.26\linewidth}>{\rcc}p{0.1\linewidth}@{}}
\toprule
\textbf{Wait} & \textbf{Applies to} & \textbf{Default} & \textbf{Adaptive} \\
\midrule
Down wait $k_{\downarrow}$ & \textsc{Anchored}$\to$\textsc{Suspicious}, \textsc{Recovering}$\to$\textsc{Sovereign} & 2, doubles per flap level up to 8 & Yes \\
Suspicious exit $k_{\uparrow}$ & \textsc{Suspicious}$\to$\textsc{Anchored} & 2 & No \\
Recovery $H$ & \textsc{Recovering}$\to$\textsc{Anchored} & 2 & No \\
\bottomrule
\end{tabular}
\end{table}

\subsection{Transitions and Hysteresis}
\label{subsec:transitions}

A \emph{warning} (suspicious Bitcoin gap, unhealthy DA, or unhealthy P2P
quality) demotes \textsc{Anchored} to \textsc{Suspicious} only after
$k_{\downarrow}$ consecutive non-healthy blocks. A \emph{critical} condition
(excessive Bitcoin gap, failed anchor re-verification, complete anchor-peer
isolation, or a \textsc{Suspicious} timeout of $T_{\text{sus}}$ blocks) moves
the chain to \textsc{Sovereign} immediately and bypasses every counter.
Recovery is deliberately slower than degradation:
\textsc{Suspicious}$\to$\textsc{Anchored} requires $k_{\uparrow}$ consecutive
healthy blocks, and \textsc{Recovering}$\to$\textsc{Anchored} requires the
safe-block counter to reach $H$ plus the exit gate of
Section~\ref{subsec:reanchoring}. \textsc{Sovereign}$\to$\textsc{Recovering}
needs only a healthy certificate.

To resist an adversary who repeatedly toggles a peripheral, the down
threshold $k_{\downarrow}$ is adaptive. It starts at a genesis value and
doubles per flap level, up to a genesis cap, and both demotion edges
(\textsc{Anchored}$\to$\textsc{Suspicious} and the non-critical
\textsc{Recovering}$\to$\textsc{Sovereign}) use it. The flap level rises by
one for each such demotion until the cap is reached, falls by one after
$T_{\text{stable}}$ healthy \textsc{Anchored} blocks, and is not reset by a
successful recovery.

The counters follow the committed state. The unhealthy streak counts
non-healthy blocks absorbed in \textsc{Anchored} or \textsc{Recovering}; it
resets on any transition and on a healthy block. The safe-block counter in
\textsc{Recovering} rises on healthy blocks up to $H$ and falls by one on an
absorbed non-healthy block; the corresponding \textsc{Suspicious} counter
behaves the same way up to $k_{\uparrow}$. Any other transition resets both.
Section~\ref{subsec:rq3} evaluates all three hysteresis-gated edges.

\subsection{Extended Proposal and Validity Predicate}
\label{subsec:proposal}

A proposal $P_h$ at height $h$ carries, in $\mathit{Txs}[0]$, the certificate
$C^{h}$, the target state $s^{h}$, a health flag, a DA receipt (last
DA-verified height and attestation flag), a Bitcoin receipt (checkpoint
height and block hash), and, when applicable, a recovery-proof reference.

\refstepcounter{definition}\label{def:validity}\textbf{Definition~\thedefinition{} (Engram proposal validity).} \texttt{ProcessProposal}
accepts a proposal $P_h$ iff
\begin{equation}
\mathsf{Valid}_h(P_h) \iff \mathsf{WellFormed}(C^{h}) \wedge \mathsf{StateOK}
\wedge \mathsf{ReceiptsOK} \wedge \mathsf{WithdrawalsOK} \wedge
\mathsf{RecoveryOK},
\label{eq:valid}
\end{equation}
where
\begin{itemize}
  \item $\mathsf{WellFormed}(C^{h})$: $C^{h}$ satisfies conditions (a)--(d) of
  Section~\ref{subsec:health-certificate};
  \item $\mathsf{StateOK}$: $s^{h} = \delta(s^{h-1}, \mathsf{Agg}(C^{h}))$ and
  the health flag matches $\mathsf{Agg}(C^{h})$;
  \item $\mathsf{ReceiptsOK}$: in \textsc{Anchored} and \textsc{Recovering},
  and for any layer the certificate reports healthy, the DA and Bitcoin
  receipts are consistent with $\mathsf{Agg}(C^{h})$ within the round-dependent
  tolerances $\tau_{\text{BTC}}$ and $\tau_{\text{DA}}$ defined below;
  \item $\mathsf{WithdrawalsOK}$: no withdrawal is included while $s^{h}$ is
  \textsc{Sovereign} or \textsc{Recovering};
  \item $\mathsf{RecoveryOK}$: the proposal carries a recovery-proof reference
  with a fresh DA attestation exactly when it targets \textsc{Recovering} with
  the safe-block counter at $H$, and the exit
  \textsc{Recovering}$\to$\textsc{Anchored} is taken only on a proof the chain
  has verified (Section~\ref{subsec:reanchoring}).
\end{itemize}
\texttt{ProcessProposal} never reads a live sensor. The FSM state is written
to replicated state only by \texttt{PreBlocker} after commit, from the
certificate already agreed upon.

The tolerances are deterministic functions of the round $r$ and genesis
parameters, so every validator computes the same bounds. Each is a base
threshold plus a widening $w_r$ that is capped, so no number of round skips
enlarges it beyond its maximum. A DA receipt may lag the Engram height by
$\tau_{\text{DA}} = \Theta_{\text{DA}} + w^{\text{DA}}_r$ Engram blocks, where
$\Theta_{\text{DA}}$ is the DA threshold and $w^{\text{DA}}_r = 0$ for
$r \le 1$, $2$ for $r = 2$, and $4$ for $r \ge 3$. A Bitcoin receipt's
checkpoint height may lag the Bitcoin tip by
$\tau_{\text{BTC}} = K_{\text{DEEP\_FINALITY}} + 2 + w^{\text{BTC}}_r$ Bitcoin
blocks, where $w^{\text{BTC}}_r = 0$ for $r \le 2$ and $1$ for $r \ge 3$. The
base $K_{\text{DEEP\_FINALITY}} + 2$ is the depth a confirmed checkpoint
already lags the tip plus two blocks of propagation slack.
Theorem~\ref{thm:withdrawal} takes both tolerances at the round in which
$B_h$ committed.

$\mathsf{ReceiptsOK}$ is exactly what a \emph{hard-gated} modular chain
enforces unconditionally. Engram's only liveness-relevant change is to make
it conditional on the certified state: when the certificate reports a layer
unhealthy outside \textsc{Anchored} and \textsc{Recovering}, that layer's
receipt is no longer required, and the circuit breaker
($\mathsf{WithdrawalsOK}$) removes the operations whose safety depended on
it.

\subsection{Recovery Proof and Re-anchoring}
\label{subsec:reanchoring}

\textbf{Normative protocol.} The exit defined here assumes proof-only recovery:
it requires a recovery proof that the chain has verified. The prototype also
contains an optional engineering heuristic that can satisfy the proof
condition; Section~\ref{subsec:limitations} identifies it.

\textbf{Statement.} On leaving \textsc{Sovereign}, a prover builds a
Noir circuit proof (Barretenberg UltraHonk
backend~\cite{aztec_noir_barretenberg,gabizon2019plonk}) over a contiguous
range of $n \leq N_{\max} = 256$ headers starting from the last anchored root
$\mathit{rt}_{\text{last}}$. A header is the tuple $(\mathit{prev\_hash},
\mathit{fsm\_state}, \mathit{withdrawal\_locked}, \mathit{state\_root})$,
where $\mathit{state\_root}$ is CometBFT's application hash for that block.
The circuit checks that (a) each header's $\mathit{prev\_hash}$ equals the
Poseidon2 hash of its predecessor, the first links to $\mathit{rt}_{\text{last}}$,
and the hash of the last header equals the claimed $\mathit{rt}_{\text{new}}$;
(b) every header has $\mathit{fsm\_state}$ in \{\textsc{Sovereign},
\textsc{Recovering}\} and $\mathit{withdrawal\_locked}$ set; and (c) the first
and last blocks of the range carry ed25519 precommit signatures (round 0) from
at least two thirds of the voting power of a validator set whose hash is a
public input (the compiled circuit fixes the set size at seven). Unused
slots are padding, so proving time, verification time, and proof size are
constant in $n$ (Appendix~\ref{app:backend}). The circuit does not check each
transition against $\delta$, and the lock flag is the one the chain recorded
in the header.

\textbf{Intended verifier.} A CometBFT quorum certificate shows that a quorum
signed a block. The recovery proof adds a statement about the interval
between two signed blocks: every recorded state in it is \textsc{Sovereign}
or \textsc{Recovering} with the lock flag set. This matters to a verifier that
does not replay the chain, such as a bridge or light client on another chain,
a node syncing from genesis, or a Bitcoin observer
(Section~\ref{subsec:security-semantics}). For validators, which already hold
the headers, it is a consistency check. The target verifier is an external
light client that accepts a re-anchoring checkpoint only with a valid proof;
the prototype does not include such a client.

\textbf{Longer intervals.} An interval longer than 256 headers is covered by
a sequence of proofs: each accepted proof advances the on-chain
$\mathit{rt}_{\text{last}}$ to its $\mathit{rt}_{\text{new}}$, and the next
proof must start from that root. This costs $\lceil L/256 \rceil$ proofs for
$L$ blocks. An aggregator circuit can fold up to four chunk proofs into one
proof of the same size, which makes verification cost independent of $L$ up to
1024 headers at the price of one extra proof
(Appendix~\ref{app:aggregation}).

\textbf{Who verifies what.} The recovery proof is submitted as a message and
verified by \emph{every Engram validator} during block execution, using an
embedded, deterministic WebAssembly build of the UltraHonk verifier. The
message handler binds the proof's public inputs to chain state:
\begin{itemize}
  \item $\mathit{rt}_{\text{last}}$ to the committed last anchored root, and
  $\mathit{rt}_{\text{new}}$ to the accumulator of a header the chain recorded;
  \item the first height to the last anchored height plus one, and the header
  count to the distance between the two checkpoints;
  \item the first and last block hashes to the CometBFT hashes the chain
  recorded at those heights, and the validator-set hash to that of the genesis
  set.
\end{itemize}
A verified height is latched so a proof cannot be replayed. Interior headers
are bound only through the accumulator, not through per-block signatures.
Bitcoin only timestamps the checkpoint; it does not verify the proof. After
recovery, validators resume writing \texttt{OP\_RETURN} checkpoints of 44 bytes
(a 4-byte tag, the Engram height, and the CometBFT block hash at that height),
and the Bitcoin gap sensor requires them to reach depth
$K_{\text{DEEP\_FINALITY}}$. Anchor re-verification compares the hash with the
local block hash. The proof's $\mathit{rt}_{\text{new}}$ is a Poseidon2
accumulator, not a block hash, and is not written to Bitcoin.

\textbf{Proof timing.} The proof message executes in block $h$ and latches the
height of the new checkpoint. The \texttt{PreBlocker} of block $h{+}1$
recomputes the proof-validity flag from that height, and
\texttt{ProcessProposal} reads this committed flag when it evaluates the exit
edge, so \textsc{Recovering}$\to$\textsc{Anchored} cannot occur before block
$h{+}2$.
\texttt{ProcessProposal} does not verify the proof itself.

\textbf{Resulting guarantee.} The recovery proof guarantees, to anyone who
verifies it, that (i) the recorded headers of the \textsc{Sovereign} interval
form a single hash chain between two boundary blocks signed by at least two
thirds of the voting power, and (ii) every recorded state is
\textsc{Sovereign} or \textsc{Recovering} with the lock flag set. Because the
lock flag is the chain's own record, (ii) does not prove that no withdrawal
executed. The proof does not make the interval's \emph{ordering} secure
against a Byzantine supermajority: during \textsc{Sovereign}, safety rests on
the validator quorum alone, which is precisely why withdrawals are locked.
Re-anchoring restores Bitcoin timestamping from the recovery point forward.

\textbf{Data availability of the \textsc{Sovereign} interval.} A recovery
proof shows the interval is recorded consistently, not that its data is
\emph{available}. The exit condition is therefore stronger than the proof:
\textsc{Recovering}$\to$\textsc{Anchored} requires (1) a recovery proof that
the chain has verified and (2) a DA backfill certificate -- a statement,
backed by the quorum of signed observations in the health certificate, that
the block data for every height of the \textsc{Sovereign} interval has been
published to the DA layer and is retrievable. During \textsc{Sovereign}, each
validator keeps the interval's blocks in its own store (they are ordinary
committed CometBFT blocks) and queues them for publication; on recovery the
queue is drained in height order. Each validator reports a contiguous height:
the highest height such that every height up to it has been confirmed
retrievable. The certificate takes the weighted median $\widehat d$ of these
reports. Let $e$ be the height at which \textsc{Recovering} was entered, and
let $\mathrm{safe}^{h-1}$ and $\mathrm{proof}^{h-1}$ be the committed
safe-block counter and proof-validity flag. The exit edge is taken at height
$h$ iff
\begin{equation}
\mathrm{healthy}\bigl(\mathsf{Agg}(C^{h})\bigr) \;\wedge\;
\mathrm{safe}^{h-1} = H \;\wedge\; \mathrm{proof}^{h-1} \;\wedge\;
\widehat d \,\ge\, e - 1 .
\label{eq:exit}
\end{equation}
The DA gap sensor keeps measuring head liveness, so the DA layer stays
healthy while the backfill drains. The specification abstracts the backfill
condition as a Boolean, \texttt{da\_backfill\_valid}. The prototype implements
the queue and the gate; Section~\ref{subsec:sim-to-real} describes the
remaining live validation.
